\section{Through the Weeks}

\subsection{The Hour to Rise from Sleep}\label{sec:first-week}

Advent begins on a Saturday evening. The Office of the First Sunday opens at
Vespers the evening before, its First Vespers, and the short reading sung there
is from St Paul to the Romans. The next morning the Missal heads the Mass with
the name of a Roman church, its station: \latin{Statio ad S. Mariam Majorem},
at St Mary Major, in Rome. The season sets out from the church of the Mother
of God. Remember the name.

The Mass opens with \latin{Ad te levavi}, the words with which the
twenty-fourth psalm begins: ``To thee, O Lord, have I lifted up my soul.'' The
Epistle is from the chapter of Romans heard the evening before:

\begin{quote}
And that, knowing the season, that it is now the hour for us to rise from
sleep. For now our salvation is nearer than when we believed. The night is
passed and the day is at hand. Let us, therefore cast off the works of darkness
and put on the armour of light.\par
{\smallskip\noindent\small ---Romans 13:11--12\par}
\end{quote}

The next two verses tell how: ``not in rioting and drunkenness, not in chambering
and impurities, not in contention and envy. But put ye on the Lord Jesus
Christ''. These are the words at which St Augustine was converted. He tells the
story in his \textit{Confessions}:

\begin{quote}
So quickly I returned to the place where Alypius was sitting; for there had I
put down the volume of the apostles, when I rose thence. I grasped, opened, and
in silence read that paragraph on which my eyes first fell,---``Not in rioting
and drunkenness, not in chambering and wantonness, not in strife and envying;
but put ye on the Lord Jesus Christ, and make not provision for the flesh, to
fulfil the lusts thereof.''\par
{\smallskip\noindent\small ---St Augustine, Confessions, Book VIII\par}
\end{quote}

He read no further, he says; he did not need to. The Church sets the same lines
at the door of Advent, for every Christian who has been putting God off.

\subsubsection*{The world grown old}

The Gospel of the day is the Lord's prophecy of the end. St Gregory the Great
preached on it in St Peter's to a Roman people who had seen scourge after
scourge, and at Matins the Church reads the opening of his homily:

\begin{quote}
Our Lord and Saviour wisheth to find us ready at His second coming. Therefore He
telleth us what will be the evils of the world as it groweth old, that He may
wean our hearts from worldly affections. Here we read what great convulsions
will go before the end, that, if we will not fear God in our prosperity, we may
at least be scourged into fearing His judgment when it is at hand.\par
{\smallskip\noindent\small ---St Gregory the Great, Homily 1 on the Gospels\par}
\end{quote}

Further on he tells his people to look about them: what they used to hear of,
they now see, new evils every day, and so few of them left of a people once
without number. Then he draws the world as an old man. In youth the chest is
strong and the neck firm; in age the frame stoops, the withered neck droops,
the chest heaves with sighs, and the breath breaks off the words as he speaks.
Gregory wants to loosen the hands that cling to a world already slipping from
them. Against such a world the Church's first prayer of the season asks God
himself to come:

\begin{studyframe}
\latin{Excita, quaesumus, Domine, potentiam tuam, et veni: ut ab imminentibus
peccatorum nostrorum periculis, te mereamur protegente eripi, te liberante
salvari: Qui vivis et regnas cum Deo Patre in unitate Spiritus Sancti, Deus:
per omnia saecula saeculorum. R. Amen.}
\medskip

Bestir, O Lord, Thy might, we pray Thee and come; That, defended by Thee, we may
deserve rescue from approaching dangers brought on by our sins, and being set
free by Thee, obtain our salvation. Who livest.\par
\end{studyframe}

\latin{Excita}, ``stir up'', is the first word of four prayers in the season.
The Church says it today of God's power, next Sunday of our hearts, and again
on the Ember Friday and on the Fourth Sunday. It is a prayer for the morning,
said before the day has a chance to speak.

At Matins this week the Church begins the book of Isaiah, whose first chapters
charge a people that has forgotten God and promise that ``if your sins be as
scarlet, they shall be made as white as snow''; Appendix~\ref{sec:daily} sets
them out day by day. From the feast of St Andrew on 30 November until Christmas
there is also a short prayer to the hour of Christ's birth, to be said fifteen
times a day; it is in Appendix~\ref{sec:prayers}, with the note printed beside
it. And look at the hour: ``it is now the hour for us to rise from sleep''.
Whatever has been put off for years can be begun now.
